\scanpage{145}
% archive foliation: 144
Notons que l'image de $\eta$ dans $H^1(S,J)$, par fonctorialité, n'est autre que la classe de $\uline{\mathrm{Pic}}^1_{X/S}$ (considéré comme torseur sous $J$). Supposons cette classe nulle [ce qui est le cas si $X/S$ admet une section]. \deleted{p. ex.} Alors la suite exacte
\diagram{\begin{tikzcd}
H^0(S,J)\arrow[r]&H^1(S,T)\arrow[r,dash]&H^1(S,E)\arrow[r,dash]&H^1(S,J)
\end{tikzcd}}
montre que \deleted{$\unsure{\eta}$} $\eta$ est dans $\operatorname{Im}H^1(S,T)$.

Notons qu'on a d'ailleurs une suite exacte \inserted{(si $Y\mathbin{\text{---}}S$ surjectif)}
\diagram{(*)\quad\begin{tikzcd}
0\arrow[r]&G_{mS}\arrow[r]&\prod_{Y/S}G_m\arrow[r]&T\arrow[r,dash]&0
\end{tikzcd}}
qui donne une suite exacte \deleted{\unsure{longue}}
\diagram{\begin{tikzcd}[column sep=small]
\unsure{0}\arrow[r]&H^0(S,O_S)^*\arrow[r,dash]&H^0(Y,O_Y)^*\arrow[r,dash]&H^0(S,T)\arrow[r,dash]&\mathrm{Pic}(S)\arrow[r]&\mathrm{Pic}(Y)\arrow[r,dash]&H^1(S,T)\arrow[r]&\mathrm{Br}'(S)\arrow[r]&\mathrm{Br}'(Y),
\end{tikzcd}}

Lorsque la suite exacte est splittée, p. ex. lorsque $Y$ admet une section sur $S$, alors on trouve
\[
\deleted{H^1(S,T)\unsure{\simeq}}\qquad H^1(S,T)\simeq\mathrm{Pic}(Y)/\operatorname{Im}\mathrm{Pic}(S)
\]
et lorsque $Y=\coprod_{i\in I}g_i(S)$ comme dessus, alors
\[
\deleted{\unsure{on trouve}}\qquad\mathrm{Pic}(Y)\simeq\mathrm{Pic}(S)^I
\]
et on trouve
\[
H^1(S,T)\simeq\mathrm{Pic}(S)^I/\operatorname{Im}\mathrm{Pic}(\unsure{S})
\]
\inserted{$\longrightarrow$ par diagonale.}

Si p. ex. $\mathrm{Pic}(S)=0$, on trouve $\eta=0$, à fortiori les $\xi_n$, $\xi_\ell$ sont nuls.

